\section{Policy Inputs and the Return to CLEWS}
\label{sec:illustration}

\subsection{Carbon price}

An exogenous carbon-price path can enter CLEWS as an emissions penalty. It can also enter
OG-Core as a payment on household energy use when that expenditure has not already been
priced through the electricity-cost ratio. Revenue can then be returned through household
transfers or used elsewhere in the public budget.

The same carbon cost cannot enter twice. If the carbon penalty has already changed the
CLEWS generation plan and is embodied in the electricity-cost ratio, OG-Core does not add
it again to the same electricity expenditure. Industrial carbon is represented through
CLEWS, because firms do not yet purchase an explicit energy input in OG-Core.

The carbon-price mapping has been tested on the economic side, and CLEWS can receive the
penalty path. The reported one-pass results do not include a CLEWS re-solve under that
policy path, so the carbon price has not yet changed the reported technology mix. Its
monetary magnitude also depends on the unfinished unit conversion.

\subsection{Private investment incentive}

The link can represent support for private generation investment through an investment
tax credit or related change in business taxation. This lowers the electricity industry's
cost of capital and shifts its demand for private capital. It is economically distinct
from the generation capital-intensity relationship, which changes production shares but
does not subsidise investment.

The mapping is defined but has not been exercised in the current validation runs. Its use will
also require a calibrated monetary bridge and a policy specification that distinguishes
the incentive from the underlying generation investment.

\subsection{Equilibrium return}

After OG-Core solves, the link writes its equilibrium portfolio return as a candidate
CLEWS discount rate. This can align resource planning with the economy's opportunity cost
of capital: a lower rate gives more weight to the future operating savings of
capital-intensive, long-lived assets.

The two rates are not automatically equivalent. A planning rate may reflect social
objectives, while an equilibrium return reflects market conditions. Risk, tax treatment,
maturity, and real-versus-nominal conventions also differ. The link therefore exposes a
choice that analysts must harmonise; it does not determine the correct social discount
rate. The return path is emitted but has not yet been passed through a repeated solution.

\subsection{Economic activity and electricity demand}

OG-Core also writes changes in industry output or household consumption to a revised
CLEWS demand path. The relevant measure depends on the service: industrial electricity
should follow the industries that use it, while household demand should follow household
consumption rather than aggregate GDP.

This return path supplies the demand response that is absent within the current CLEWS
optimisation. Higher electricity costs may reduce use and therefore change the capacity
required. Investment and health effects may raise activity and demand in other periods.

The mechanism for copying a revised demand path into CLEWS, solving the model, and reading
the result back has been tested with a controlled demand change. The activity path emitted
by OG-Core has not yet travelled through the complete outer loop. Convergence and the
propagation of uncertainty across repeated solutions therefore remain untested.
